\editionPage[\pdfbookmark[0]{Deux bilans annuels}{bm-bilans-guide-parent}]{bilans-guide}{Deux bilans pour l’année}{Sujets originaux · Relier les chapitres}{Durée conseillée : 1 h 15 · Barème d’entraînement sur 20}
\begin{blocEdition}{bilans-guide}

Ces évaluations sont des créations du livre pour le CM1. Elles n’ont pas
de statut d’examen officiel. Le temps et le barème sont des repères pour
une séance autonome ; ils s’adaptent librement à la classe ou à la maison.
\editionEtape{Automatismes : environ 15 minutes}
Douze questions courtes, sans calculatrice ni brouillon long : 6 points,
soit 0,5 point par réponse. Écrire les étapes utiles, même brèves.
\editionEtape{Problèmes : environ 1 heure}
Quatre problèmes indépendants : 14 points. Chaque problème vaut
3,5 points : trois questions à 1 point et 0,5 point pour la clarté de la
rédaction, les unités et les contrôles.
\editionEtape{Travailler après l’évaluation}
Comparer les démarches, relever les erreurs par notion, puis refaire
les questions utiles sans correction. Une démarche peut être juste
même quand le résultat final est faux : repérer l’étape à reprendre.
\editionEtape{Préparer une autre séance}
Revenir au cours et à deux exercices du chapitre concerné. Le second
bilan peut être fait plus tard pour vérifier que la méthode est acquise.

\end{blocEdition}
\editionPage{bilan-A-auto}{Bilan A}{Nombres, calculs et premiers outils}{Partie 1 · 15 min · Sans calculatrice · 6 points}
\begin{blocEdition}{bilan-A-auto}
\begin{questions}
\item Écrire le nombre $8\times1\,000+3\times100+7$.
\item Calculer de tête $457+19$.
\item Poser et calculer $702-458$.
\item Calculer $36:4$.
\item Calculer $7\times60$.
\item Encadrer $3\,456$ entre deux centaines consécutives.
\item Calculer $\frac38+\frac28$.
\item Écrire avec une virgule : $2+\frac3{10}+\frac5{100}$.
\item Comparer $4{,}5$ et $4{,}45$.
\item Calculer $3{,}6\times10$.
\item Compléter : $25\times\square=100$.
\item Donner l’arrondi à l’entier de $6{,}78$.\end{questions}
\end{blocEdition}
\editionPage{bilan-A-ex-1}{Bilan A · Problème 1}{La fête de l’école}{Partie 2 · 1 heure pour les quatre problèmes · 3,5 points}
\begin{blocEdition}{bilan-A-ex-1}
\begin{questions}
\item La fête de l’école a rapporté 875 euros aux stands de jeux et 468 euros à la buvette. Calculer la recette totale.
\item Les dépenses s’élèvent à 390 euros. Quelle somme reste-t-il ?
\item Le reste est partagé également entre 7 classes. Combien reçoit chaque classe ?
\end{questions}\par\vspace{5mm}{\otLight\fontsize{8}{10}\selectfont\color{otMuted}Recherche et rédaction}\par\par\vspace{7mm}{\color{otRule!65}\hrule height .25pt}\par\vspace{7mm}{\color{otRule!65}\hrule height .25pt}\par\vspace{7mm}{\color{otRule!65}\hrule height .25pt}\par\vspace{7mm}{\color{otRule!65}\hrule height .25pt}\par\vspace{7mm}{\color{otRule!65}\hrule height .25pt}\par\vspace{7mm}{\color{otRule!65}\hrule height .25pt}
\end{blocEdition}
\editionPage{bilan-A-ex-2}{Bilan A · Problème 2}{Les cartons de la médiathèque}{Partie 2 · 1 heure pour les quatre problèmes · 3,5 points}
\begin{blocEdition}{bilan-A-ex-2}
\begin{questions}
\item La médiathèque reçoit $1\,350$ livres neufs. Elle les range dans des cartons de 8 livres. Combien de cartons complets remplit-elle ?
\item Combien de livres restent hors carton ?
\item Combien de cartons faut-il au total pour transporter tous les livres ?
\end{questions}\par\vspace{5mm}{\otLight\fontsize{8}{10}\selectfont\color{otMuted}Recherche et rédaction}\par\par\vspace{7mm}{\color{otRule!65}\hrule height .25pt}\par\vspace{7mm}{\color{otRule!65}\hrule height .25pt}\par\vspace{7mm}{\color{otRule!65}\hrule height .25pt}\par\vspace{7mm}{\color{otRule!65}\hrule height .25pt}\par\vspace{7mm}{\color{otRule!65}\hrule height .25pt}\par\vspace{7mm}{\color{otRule!65}\hrule height .25pt}
\end{blocEdition}
\editionPage{bilan-A-ex-3}{Bilan A · Problème 3}{Le gâteau du dimanche}{Partie 2 · 1 heure pour les quatre problèmes · 3,5 points}
\begin{blocEdition}{bilan-A-ex-3}
\begin{questions}
\item Un gâteau est coupé en 8 parts égales. La famille en mange 5 parts le midi. Quelle fraction du gâteau reste-t-il ?
\item Le soir, chacun des 3 convives mange une part. Quelle fraction du gâteau a été mangée sur la journée ?
\item Un second gâteau identique est partagé entre 4 enfants. Quelle est la part de chacun ? La donner en fraction de gâteau.
\end{questions}\par\vspace{5mm}{\otLight\fontsize{8}{10}\selectfont\color{otMuted}Recherche et rédaction}\par\par\vspace{7mm}{\color{otRule!65}\hrule height .25pt}\par\vspace{7mm}{\color{otRule!65}\hrule height .25pt}\par\vspace{7mm}{\color{otRule!65}\hrule height .25pt}\par\vspace{7mm}{\color{otRule!65}\hrule height .25pt}\par\vspace{7mm}{\color{otRule!65}\hrule height .25pt}\par\vspace{7mm}{\color{otRule!65}\hrule height .25pt}
\end{blocEdition}
\editionPage{bilan-A-ex-4}{Bilan A · Problème 4}{Les courses du mercredi}{Partie 2 · 1 heure pour les quatre problèmes · 3,5 points}
\begin{blocEdition}{bilan-A-ex-4}
\begin{questions}
\item Lou achète un livre à $8{,}70$ euros et un jeu de cartes à $4{,}55$ euros. Calculer la dépense.
\item Elle paie avec un billet de 20 euros. Combien lui rend-on ?
\item Avec la monnaie rendue, peut-elle acheter trois carnets à $2{,}20$ euros l’un ?
\end{questions}\par\vspace{5mm}{\otLight\fontsize{8}{10}\selectfont\color{otMuted}Recherche et rédaction}\par\par\vspace{7mm}{\color{otRule!65}\hrule height .25pt}\par\vspace{7mm}{\color{otRule!65}\hrule height .25pt}\par\vspace{7mm}{\color{otRule!65}\hrule height .25pt}\par\vspace{7mm}{\color{otRule!65}\hrule height .25pt}\par\vspace{7mm}{\color{otRule!65}\hrule height .25pt}\par\vspace{7mm}{\color{otRule!65}\hrule height .25pt}
\end{blocEdition}
\editionPage{bilan-B-auto}{Bilan B}{Toute l’année en quatre problèmes}{Partie 1 · 15 min · Sans calculatrice · 6 points}
\begin{blocEdition}{bilan-B-auto}
\begin{questions}
\item Donner le chiffre des dizaines de mille de $582\,741$.
\item Calculer de tête $68\times5$.
\item Calculer de tête $801-599$.
\item Donner le quotient et le reste de $75:8$.
\item Écrire $\frac74$ avec un entier et une fraction.
\item Écrire $0{,}75$ sous forme de fraction.
\item Comparer $5{,}08$ et $5{,}8$.
\item Calculer $10-3{,}45$.
\item Convertir 2 h 15 min en minutes.
\item Donner le périmètre d’un carré de côté 9 cm.
\item Comparer 3 kg et $2\,900$ g.
\item Quatre pains coûtent 6 euros. Combien coûtent 8 pains ?\end{questions}
\end{blocEdition}
\editionPage{bilan-B-ex-1}{Bilan B · Problème 1}{Le train de Lucie}{Partie 2 · 1 heure pour les quatre problèmes · 3,5 points}
\begin{blocEdition}{bilan-B-ex-1}
\begin{questions}
\item Lucie prend un train qui part à 9 h 47 et arrive à 12 h 15. Combien dure le voyage ?
\item Son retour dure 2 h 40 et se termine à 19 h 05. À quelle heure est-elle partie ?
\item Quel trajet est le plus long, et de combien de minutes ?
\end{questions}\par\vspace{5mm}{\otLight\fontsize{8}{10}\selectfont\color{otMuted}Recherche et rédaction}\par\par\vspace{7mm}{\color{otRule!65}\hrule height .25pt}\par\vspace{7mm}{\color{otRule!65}\hrule height .25pt}\par\vspace{7mm}{\color{otRule!65}\hrule height .25pt}\par\vspace{7mm}{\color{otRule!65}\hrule height .25pt}\par\vspace{7mm}{\color{otRule!65}\hrule height .25pt}\par\vspace{7mm}{\color{otRule!65}\hrule height .25pt}
\end{blocEdition}
\editionPage{bilan-B-ex-2}{Bilan B · Problème 2}{Le potager de l’école}{Partie 2 · 1 heure pour les quatre problèmes · 3,5 points}
\begin{blocEdition}{bilan-B-ex-2}
\begin{questions}
\item Le potager est un rectangle de $7{,}5$ m sur 4 m. Calculer son périmètre.
\item On le clôture avec un grillage vendu $2{,}50$ euros le mètre. Calculer le prix du grillage.
\item La classe dispose de 60 euros. Quelle somme restera-t-il après l’achat ?
\end{questions}\par\vspace{5mm}{\otLight\fontsize{8}{10}\selectfont\color{otMuted}Recherche et rédaction}\par\par\vspace{7mm}{\color{otRule!65}\hrule height .25pt}\par\vspace{7mm}{\color{otRule!65}\hrule height .25pt}\par\vspace{7mm}{\color{otRule!65}\hrule height .25pt}\par\vspace{7mm}{\color{otRule!65}\hrule height .25pt}\par\vspace{7mm}{\color{otRule!65}\hrule height .25pt}\par\vspace{7mm}{\color{otRule!65}\hrule height .25pt}
\end{blocEdition}
\editionPage{bilan-B-ex-3}{Bilan B · Problème 3}{L’enquête des déplacements}{Partie 2 · 1 heure pour les quatre problèmes · 3,5 points}
\begin{blocEdition}{bilan-B-ex-3}
\begin{questions}
\item Une enquête interroge les 55 élèves de deux classes.\begin{center}\renewcommand{\arraystretch}{1.35}\begin{tabular}{c|ccc}\toprule
 & À pied & En bus & À vélo\\\midrule
CM1 & 12 & 9 & 5\\
CM2 & 10 & 13 & 6\\
\bottomrule\end{tabular}\end{center}Combien d’élèves viennent en bus ?
\item Dans quelle classe la marche à pied est-elle la plus fréquente ?
\item Vérifier que le tableau recense bien les 55 élèves.
\end{questions}\par\vspace{5mm}{\otLight\fontsize{8}{10}\selectfont\color{otMuted}Recherche et rédaction}\par\par\vspace{7mm}{\color{otRule!65}\hrule height .25pt}\par\vspace{7mm}{\color{otRule!65}\hrule height .25pt}\par\vspace{7mm}{\color{otRule!65}\hrule height .25pt}\par\vspace{7mm}{\color{otRule!65}\hrule height .25pt}\par\vspace{7mm}{\color{otRule!65}\hrule height .25pt}\par\vspace{7mm}{\color{otRule!65}\hrule height .25pt}
\end{blocEdition}
\editionPage{bilan-B-ex-4}{Bilan B · Problème 4}{La kermesse du hasard}{Partie 2 · 1 heure pour les quatre problèmes · 3,5 points}
\begin{blocEdition}{bilan-B-ex-4}
\begin{questions}
\item Au stand de la kermesse, une roue porte 10 secteurs identiques : 6 marqués « perdu », 3 « sucette », 1 « cadeau ». Quel résultat est le plus probable ?
\item A-t-on une chance sur deux d’obtenir une sucette ? Justifier.
\item Le stand propose aussi 2 tailles de ballons et 5 couleurs. Combien de ballons différents peut-on composer ?
\end{questions}\par\vspace{5mm}{\otLight\fontsize{8}{10}\selectfont\color{otMuted}Recherche et rédaction}\par\par\vspace{7mm}{\color{otRule!65}\hrule height .25pt}\par\vspace{7mm}{\color{otRule!65}\hrule height .25pt}\par\vspace{7mm}{\color{otRule!65}\hrule height .25pt}\par\vspace{7mm}{\color{otRule!65}\hrule height .25pt}\par\vspace{7mm}{\color{otRule!65}\hrule height .25pt}\par\vspace{7mm}{\color{otRule!65}\hrule height .25pt}
\end{blocEdition}
